% GENERATED FILE. Edit canonical lessons, then run npm run generate:books.
% canonical-source-hash: fnv1a64:7c2d8cd10a0d8761
% canonical-lessons: RU-C471-obida, RU-C471-obizhatsya, RU-C471-obizhat, RU-C471-obidchivyy, RU-C471-proshchat

\chapter{Hurt Feelings}
\label{ch:ru-hurt-feelings}
\hlchaptermodality{471}

\begin{chapteropening}
\textbf{By the end of this chapter:} \emph{I can talk about hurt feelings and forgiving.}

Complete the last lesson of chapter 471: I can talk about hurt feelings and forgiving.
\end{chapteropening}

\section[\texorpdfstring{\ru{обида} (obída)}{обида (obída)}]{\ru{обида} (obída) — hurt feelings, offence}

\label{lesson:RU-C471-obida}



\begin{quote}
Before the new one: say the Russian for to blush, then the Russian for shy.
\end{quote}

\subsection*{You'll want to know: \ru{обида}}
\textbf{\ru{обида}} (\emph{obída}) — \textquotedblleft{}hurt feelings, offence\textquotedblright{}.

\begin{quote}
\ru{Он} \ru{до} \ru{сих} \ru{пор} \ru{помнит} \ru{эту} \ru{обиду}.
\end{quote}

\begin{quote}
\emph{On do sikh pór pómnit étu obídu.}
\end{quote}

\begin{quote}
He still remembers that hurt.
\end{quote}

Feminine. \emph{\ru{Не} \ru{в} \ru{обиду}} — no offence meant.

\begin{cousinweb}
\emph{\ru{обида}} is inherited from Proto-Slavic \emph{obvida} (EtymDB), derived from \emph{obviděti}, which is built on \emph{ob-} and \emph{viděti} 'to see'.
\end{cousinweb}

\begin{grammarlens}[title={What you've built}]
One more word for telling, explaining and describing at B1.
\end{grammarlens}

\subsection*{Your turn}
\begin{itemize}
\raggedright
  \item \emph{Say it:} \emph{obída}
  \item \emph{Say it:} \emph{obída}, once more
  \item \emph{Say it:} a sentence of your own with \emph{obída}
  \item \emph{Recall:} say the Russian for to blush, then the Russian for shy, then say \emph{obída} again
\end{itemize}

\begin{tcolorbox}[breakable,colback=teal!4,colframe=teal!35!black,title={Before you move on}]
What does \ru{обида} mean? (\textquotedblleft{}Hurt feelings, offence\textquotedblright{}.) Say it once more.
\end{tcolorbox}
\section[\texorpdfstring{\ru{обижаться} (obizhátsya)}{обижаться (obizhátsya)}]{\ru{обижаться} (obizhátsya) — to take offence}

\label{lesson:RU-C471-obizhatsya}



\begin{quote}
Before the new one: say the Russian for shy, then the Russian for hurt feelings.
\end{quote}

\subsection*{You'll want to know: \ru{обижаться}}
\textbf{\ru{обижаться}} (\emph{obizhátsya}) — \textquotedblleft{}to take offence\textquotedblright{}.

\begin{quote}
\ru{Она} \ru{часто} \ru{обижается} \ru{на} \ru{мои} \ru{шутки}.
\end{quote}

\begin{quote}
\emph{Oná chásto obizháyetsya na moí shútki.}
\end{quote}

\begin{quote}
She often takes offence at my jokes.
\end{quote}

Imperfective; the perfective is \ru{обидеться} (\ru{обижусь}, \ru{обидишься}). \emph{\ru{Обижаться} \ru{на} \ru{кого}} — to be offended with someone.

\begin{cousinweb}
\emph{\ru{обижаться}} is the reflexive of \emph{\ru{обижать}}, the imperfective partner of \emph{\ru{обидеть}}; all three share the root of \emph{\ru{обида}} 'offence', from Proto-Slavic \emph{obvida} (EtymDB).
\end{cousinweb}

\begin{grammarlens}[title={What you've built}]
One more word for telling, explaining and describing at B1.
\end{grammarlens}

\subsection*{Your turn}
\begin{itemize}
\raggedright
  \item \emph{Say it:} \emph{obizhátsya}
  \item \emph{Say it:} \emph{obizhátsya}, once more
  \item \emph{Say it:} a sentence of your own with \emph{obizhátsya}
  \item \emph{Recall:} say the Russian for shy, then the Russian for hurt feelings, then say \emph{obizhátsya} again
\end{itemize}

\begin{tcolorbox}[breakable,colback=teal!4,colframe=teal!35!black,title={Before you move on}]
What does \ru{обижаться} mean? (\textquotedblleft{}To take offence\textquotedblright{}.) Say it once more.
\end{tcolorbox}
\section[\texorpdfstring{\ru{обижать} (obizhát)}{обижать (obizhát)}]{\ru{обижать} (obizhát) — to hurt (someone's feelings)}

\label{lesson:RU-C471-obizhat}



\begin{quote}
Before the new one: say the Russian for hurt feelings, then the Russian for to take offence.
\end{quote}

\subsection*{You'll want to know: \ru{обижать}}
\textbf{\ru{обижать}} (\emph{obizhát}) — \textquotedblleft{}to hurt (someone's feelings)\textquotedblright{}.

\begin{quote}
\ru{Старший} \ru{брат} \ru{никогда} \ru{не} \ru{обижал} \ru{младшего}.
\end{quote}

\begin{quote}
\emph{Stárshiy brat nikogdá ne obizhál mládshego.}
\end{quote}

\begin{quote}
The older brother never picked on the younger one.
\end{quote}

Imperfective; perfective \ru{обидеть}. Takes the accusative: \emph{\ru{обижать} \ru{младшего}}.

\begin{cousinweb}
\emph{\ru{обижать}} is the imperfective partner of \emph{\ru{обидеть}} 'to offend' (EtymDB lists both without a source); they share the root of \emph{\ru{обида}}, from Proto-Slavic \emph{obvida} (EtymDB).
\end{cousinweb}

\begin{grammarlens}[title={What you've built}]
One more word for telling, explaining and describing at B1.
\end{grammarlens}

\subsection*{Your turn}
\begin{itemize}
\raggedright
  \item \emph{Say it:} \emph{obizhát}
  \item \emph{Say it:} \emph{obizhát}, once more
  \item \emph{Say it:} a sentence of your own with \emph{obizhát}
  \item \emph{Recall:} say the Russian for hurt feelings, then the Russian for to take offence, then say \emph{obizhát} again
\end{itemize}

\begin{tcolorbox}[breakable,colback=teal!4,colframe=teal!35!black,title={Before you move on}]
What does \ru{обижать} mean? (\textquotedblleft{}To hurt (someone's feelings)\textquotedblright{}.) Say it once more.
\end{tcolorbox}
\section[\texorpdfstring{\ru{обидчивый} (obídchivyy)}{обидчивый (obídchivyy)}]{\ru{обидчивый} (obídchivyy) — touchy}

\label{lesson:RU-C471-obidchivyy}



\begin{quote}
Before the new one: say the Russian for to take offence, then the Russian for to hurt.
\end{quote}

\subsection*{You'll want to know: \ru{обидчивый}}
\textbf{\ru{обидчивый}} (\emph{obídchivyy}) — \textquotedblleft{}touchy\textquotedblright{}.

\begin{quote}
\ru{С} \ru{обидчивым} \ru{человеком} \ru{трудно} \ru{шутить}.
\end{quote}

\begin{quote}
\emph{S obídchivym chelovékom trúdno shutít.}
\end{quote}

\begin{quote}
It's hard to joke with a touchy person.
\end{quote}

Easily offended. Short forms: \ru{обидчив}, \ru{обидчива}.

\begin{cousinweb}
\emph{\ru{обидчивый}} is built on \emph{\ru{обида}} 'offence' with the ending \emph{-\ru{чивый}}, which marks a habit; \emph{\ru{обида}} is inherited from Proto-Slavic \emph{obvida} (EtymDB).
\end{cousinweb}

\begin{grammarlens}[title={What you've built}]
One more word for telling, explaining and describing at B1.
\end{grammarlens}

\subsection*{Your turn}
\begin{itemize}
\raggedright
  \item \emph{Say it:} \emph{obídchivyy}
  \item \emph{Say it:} \emph{obídchivyy}, once more
  \item \emph{Say it:} a sentence of your own with \emph{obídchivyy}
  \item \emph{Recall:} say the Russian for to take offence, then the Russian for to hurt, then say \emph{obídchivyy} again
\end{itemize}

\begin{tcolorbox}[breakable,colback=teal!4,colframe=teal!35!black,title={Before you move on}]
What does \ru{обидчивый} mean? (\textquotedblleft{}Touchy\textquotedblright{}.) Say it once more.
\end{tcolorbox}
\section[\texorpdfstring{\ru{прощать} (proshchát)}{прощать (proshchát)}]{\ru{прощать} (proshchát) — to forgive}

\label{lesson:RU-C471-proshchat}



\begin{quote}
Before the new one: say the Russian for to hurt, then the Russian for touchy.
\end{quote}

\subsection*{You'll want to know: \ru{прощать}}
\textbf{\ru{прощать}} (\emph{proshchát}) — \textquotedblleft{}to forgive\textquotedblright{}.

\begin{quote}
\ru{Мама} \ru{всегда} \ru{прощает} \ru{нас}.
\end{quote}

\begin{quote}
\emph{Máma vsegdá proshcháyet nas.}
\end{quote}

\begin{quote}
Mum always forgives us.
\end{quote}

Imperfective; perfective \ru{простить} (\ru{прощу}, \ru{простишь}), the verb behind \emph{\ru{простите}}. \emph{\ru{Прощать} \ru{кому} \ru{что}}: \ru{прощать} \ru{другу} \ru{ошибку}.

\begin{cousinweb}
\emph{\ru{прощать}} is the imperfective partner of \emph{\ru{простить}} 'to forgive'; EtymDB links it to Old East Slavic \emph{prostiti}, and both verbs are usually traced to \emph{\ru{простой}} 'simple, free', so to forgive is to set someone free.
\end{cousinweb}

\begin{grammarlens}[title={What you've built}]
One more word for telling, explaining and describing at B1.
\end{grammarlens}

\subsection*{Your turn}
\begin{itemize}
\raggedright
  \item \emph{Say it:} \emph{proshchát}
  \item \emph{Say it:} \emph{proshchát}, once more
  \item \emph{Say it:} a sentence of your own with \emph{proshchát}
  \item \emph{Recall:} say the Russian for to hurt, then the Russian for touchy, then say \emph{proshchát} again
\end{itemize}

\begin{tcolorbox}[breakable,colback=teal!4,colframe=teal!35!black,title={Before you move on}]
What does \ru{прощать} mean? (\textquotedblleft{}To forgive\textquotedblright{}.) Say it once more.
\end{tcolorbox}
